%% ASPLOS 2026 project — final presentation
%% Beamer slides, ~15 min talk (35% of grade).

\documentclass[aspectratio=169,10pt]{beamer}

\usetheme{metropolis}
\usepackage{booktabs}
\usepackage{appendixnumberbeamer}
\usepackage{listings}

\lstset{
  basicstyle=\ttfamily\scriptsize,
  breaklines=true,
  columns=fullflexible,
  frame=none,
  xleftmargin=0.5em,
  numbers=none
}

\setbeamertemplate{frame footer}{Qwen2.5-Coder on Blackhole \textbullet{} MBZUAI}

\title{Optimizing Qwen2.5-Coder Decode on Tenstorrent Blackhole}
\subtitle{Two Wins, One Dead End, and What We Learned}

\author{Rassul Magauin \and Assylzhan Khamiyev \and \\
        Mohammed Rashed Ali Yahmoor Alshehhi \and \\
        Sultan Akimaliyev \and Hamad Khalifa Alyahyaee}

\institute{Mohamed bin Zayed University of Artificial Intelligence \\
           AI Systems --- Prof.\ Abdulrahman Mahmoud}

\date{April 2026}

\begin{document}

\maketitle

% ---------------------------------------------------------------------------
\begin{frame}{Where we started}
  \textbf{Direction B --- optimization.} Take Qwen2.5-Coder on
  Tenstorrent Blackhole and make it faster.
  \vfill
  \begin{itemize}
    \item Base: \texttt{tt\_transformers} already supports the Qwen2.5
          family.
    \item Hardware: 4$\times$ Blackhole p150 in tensor-parallel mesh.
    \item Method: profile, rank, fix, re-profile. Repeat three times.
  \end{itemize}
  \vfill
  \begin{alertblock}{Spoiler}
    We tried three optimizations. Two worked. One made the model
    \emph{slower}. The failure is the most interesting result.
  \end{alertblock}
\end{frame}

% ---------------------------------------------------------------------------
\section{The Machine}

\begin{frame}{Blackhole p150 in one slide}
  \begin{columns}[T,onlytextwidth]
    \begin{column}{.55\textwidth}
      \begin{itemize}
        \item 130 Tensix cores in an $8{\times}10$ compute grid
        \item Each core: FPU (matrix), SFPU (vector), 5 RISC-V CPUs,
              1.5~MB L1 SRAM
        \item 8 DRAM banks, 2D NoC
        \item 3 math-fidelity modes:
          \begin{itemize}
            \item LoFi — 4 TFLOPS, drops 4 mantissa bits
            \item HiFi2 — 2 TFLOPS, lossless on BFP8
            \item HiFi4 — 1 TFLOPS, full BF16
          \end{itemize}
      \end{itemize}
    \end{column}
    \begin{column}{.45\textwidth}
      \small
      \textbf{\texttt{tt-metal} programming model}\\[0.3em]
      Each op = a program:
      \begin{itemize}\small
        \item reader kernel pulls tiles
        \item compute kernel runs FPU/SFPU
        \item writer kernel pushes tiles
      \end{itemize}
      Circular buffers in L1 stitch them together. Fusion happens at
      the \emph{pipeline} level, not ``one big kernel''.
    \end{column}
  \end{columns}
\end{frame}

% ---------------------------------------------------------------------------
\begin{frame}{The workload}
  \textbf{Qwen2.5-Coder-7B}, architecturally identical to Qwen2.5-7B-Instruct:
  \begin{itemize}
    \item 28 decoder layers, hidden 3584, MLP intermediate 18944
    \item 28 query heads, 4 KV heads (GQA), $\mathrm{rope\_theta}=10^{6}$
    \item SwiGLU MLP: \texttt{down(silu(gate(x)) * up(x))}
  \end{itemize}
  \vfill
  \textbf{Config we profile:} \\
  batch 1, 1024-context, 10 decoder layers, paged attention, 2 generated
  tokens. Tracy profiler with
  \texttt{-a device\_kernel\_duration}.
  \vfill
  \begin{alertblock}{Why 10 layers}
    Compile time at 28 layers $\sim$6~min per run. At 10 layers we get
    comparable op-level breakdowns in under 2 minutes and can iterate.
  \end{alertblock}
\end{frame}

% ---------------------------------------------------------------------------
\section{Baseline: Where Does the Time Go?}

\begin{frame}{Baseline profile, stock recipe}
  \small
  \centering
  \begin{tabular}{lrrr}
    \toprule
    Operator / group & Time (ms) & Calls & \% \\
    \midrule
    \alert{TopK (sampling, 1-core)}    & 21.82 &   4 & \alert{52.2} \\
    \alert{Matmul (all linears)}       & 15.46 & 180 & \alert{37.0} \\
    ReduceScatter + AllGather (CCL)    &  2.01 & 123 &  4.8 \\
    BinaryNg (eltwise)                 &  0.69 & 124 &  1.7 \\
    LayerNorm / RMSNorm                &  0.37 &  63 &  0.9 \\
    SDPA decode + RoPE + reshapes      &  0.16 &  30 &  0.4 \\
    Sampling                           &  0.11 &   4 &  0.3 \\
    Other                              &  1.15 &  -  &  2.7 \\
    \midrule
    \textbf{Total}                     & \textbf{41.77} & \textbf{928} & \textbf{100.0} \\
    \bottomrule
  \end{tabular}
  \vfill
  Qwen2.5-Coder-7B decode, 4-way Blackhole TP, device~0, two trace-replay
  sessions totalled.
\end{frame}

\begin{frame}{Two things jump out}
  \begin{columns}[T,onlytextwidth]
    \begin{column}{.48\textwidth}
      \textbf{TopK at 52\%}\\[0.3em]
      Only 4 calls (one per generated token). 5.45~ms per call.\\[0.5em]
      The fused SDPA kernel on 30 calls is \textbf{under half a millisecond} total.\\[0.5em]
      A single TopK call takes 11$\times$ longer than \emph{all} attention kernels combined.
    \end{column}
    \begin{column}{.48\textwidth}
      \textbf{Matmul at 37\%}\\[0.3em]
      180 calls. MLP block (FF1+FF2+FF3) is the largest slice.\\[0.5em]
      Attention (QKV + output proj): $\sim$1.3~ms.\\[0.5em]
      MLP (gate + up + down): $\sim$3.8~ms.
    \end{column}
  \end{columns}
  \vfill
  \begin{block}{Plan}
    Fix TopK first (it is the biggest number). Then tackle the MLP.
    Leave attention alone --- the fused kernels are already doing
    their job.
  \end{block}
\end{frame}

\begin{frame}[fragile]{The TopK smell test}
\begin{lstlisting}
Op: TopKDeviceOperation
Input shape:   [1, 1, 32, 38016]
k:             32
CORE COUNT:    1     <--  one Tensix core out of 130
DEVICE KERNEL DURATION [ns]:  5,455,012
\end{lstlisting}
  \vfill
  \small
  \textbf{Input width 38016} = Qwen vocab 152{,}064 / 4 devices.
  Not a power of two.
  And \texttt{reduction/topk/} (under \texttt{ttnn/cpp/ttnn/operations/})
  has a multi-core bitonic sort implementation.\\[0.3em]
  \textbf{Hypothesis:} pad the per-device vocab up to $65{,}536 = 2^{16}$.
  Kernel flips to the multi-core path. Done.\\[0.3em]
  One-line change in \texttt{model\_config.py}. What could go wrong?
\end{frame}

% ---------------------------------------------------------------------------
\section{Attempt 1: Pad the Vocab (It Didn't Work)}

\begin{frame}{What we changed}
  In \texttt{model\_config.py}, replace\\[0.3em]
  \texttt{self.padded\_vocab\_size = 128 * 1024 if self.is\_galaxy else None}\\[0.3em]
  with a rule that rounds \emph{per-device} vocab up to the next power of two.
  For Qwen2.5-7B on 4 devices: $38016 \to 65536$.
  \vfill
  The LM head gets a wider output. TopK gets more input, but on a
  presumed-parallel kernel. Compile, re-profile, celebrate.
\end{frame}

\begin{frame}{What actually happened}
  \centering
  \small
  \begin{tabular}{lrr}
    \toprule
     & Baseline & Padded vocab \\
    \midrule
    TopK total time (ms)         & 21.82  & \alert{40.16} \\
    TopK time per call ($\mu$s)  & 5{,}455 & \alert{10{,}039} \\
    TopK CORE COUNT              &  1     & \alert{1}     \\
    Matmul total (ms)            & 15.46  & 11.84         \\
    (LM head $n$ grew)           &        &               \\
    \midrule
    \emph{Total device (ms)}     & \emph{41.77} & \alert{\emph{56.82}} \\
    \bottomrule
  \end{tabular}
  \vfill
  \textbf{TopK got 84\% slower. End-to-end 36\% slower.}
\end{frame}

\begin{frame}{Why it failed}
  The kernel \alert{never switched paths}. CORE COUNT stayed at 1.\\[0.5em]
  We just fed the same single-core kernel 1.72$\times$ more data.
  \vfill
  \begin{block}{Lesson}
    On \texttt{tt-metal}, kernel dispatch is chosen by the op's
    \emph{device operation} class based on the full tensor spec ---
    \textbf{not by input shape alone}.\\[0.5em]
    ``Round up to a friendlier shape and hope for a faster code path''
    is a GPU reflex. It does not transfer.
  \end{block}
  \vfill
  \textbf{Reverted} (commit \texttt{f6de95bb02}). Cleared the stale
  \texttt{output\_lm\_head\_*} tensor cache and re-profiled to confirm.
\end{frame}

\begin{frame}{What a real TopK fix would take}
  \begin{itemize}
    \item Parallelize \texttt{reduction::topk}'s single-core path, or
    \item Fuse TopK into \texttt{SamplingDeviceOperation}.
  \end{itemize}
  \vfill
  Both are C++ kernel work. Not a Python-level configuration tweak.
  We filed it as future work and moved on.
  \vfill
  \begin{alertblock}{Side observation}
    The stock Qwen2.5-7B recipe ships with TopK at 45--52\% of per-token
    device time. \textbf{Every user of this recipe} is paying that tax,
    not just us.
  \end{alertblock}
\end{frame}

% ---------------------------------------------------------------------------
\section{Attempt 2: MLP HiFi4 $\to$ HiFi2}

\begin{frame}[fragile]{Why HiFi2 is free here}
  \begin{itemize}
    \item HiFi4 = 1 TFLOPS on BF16, retains full mantissa
    \item HiFi2 = 2 TFLOPS, drops extra mantissa bits \emph{during the dot product}
    \item When both operands are BFP8, HiFi2 drops nothing that was there
  \end{itemize}
  \vfill
  Stock \texttt{performance\_decoder\_config.json} for Qwen2.5-7B:
\begin{lstlisting}
"LI_FF1_FF3": "HIFI4",
"LI_FF2":     "HIFI4",
\end{lstlisting}
  $\rightarrow$ change both to \texttt{"HIFI2"} in all 28 layers. 56 field edits, otherwise a no-op.
\end{frame}

\begin{frame}{Result}
  \centering
  \small
  \begin{tabular}{lrrr}
    \toprule
     & Baseline & +HiFi2 & $\Delta$ \\
    \midrule
    Matmul total (ms)              & 15.46 & \alert{10.45} & \alert{${-}32.4\%$} \\
    TopK (ms)                      & 21.82 & 21.82 & $0\%$ \\
    CCL (ms)                       &  2.01 &  2.01 & $0\%$ \\
    BinaryNg (ms)                  &  0.69 &  0.69 & $0\%$ \\
    \midrule
    \emph{Total device (ms)}       & \emph{41.77} & \alert{\emph{36.75}} & \alert{${-}12.0\%$} \\
    \bottomrule
  \end{tabular}
  \vfill
  \begin{itemize}
    \item Biggest single win in the project: \textbf{32\% off matmul time}.
    \item Accuracy: \texttt{ci-token-matching} passes. 2-layer smoke
          test decodes coherently.
  \end{itemize}
\end{frame}

% ---------------------------------------------------------------------------
\section{Attempt 3: Fuse Gate + Up Projections}

\begin{frame}{The idea}
  SwiGLU: \texttt{silu(gate\_proj(x)) * up\_proj(x)}
  \vfill
  Two matmuls share the same input $x$:
  \begin{itemize}
    \item \texttt{gate\_proj}: $[3584] \to [5120]$ per device
    \item \texttt{up\_proj}:   $[3584] \to [5120]$ per device
  \end{itemize}
  Concatenate the weights: single matmul $[3584] \to [10240]$, then slice.
  \vfill
  \textbf{tt-metal-specific reason:}
  \texttt{dram\_shard\_core\_grid\_for\_k\_and\_n} picks a different
  grid for the wider $n$. We predicted 4 cores $\to$ 8 cores $\to$
  $\approx$2$\times$ matmul speedup.
\end{frame}

\begin{frame}[fragile]{Implementation, in one place}
\begin{lstlisting}
# mlp.py __init__
self.fuse_gate_up = (not args.is_galaxy) and (prefetcher is None)
# build w_gate_up next to w1, w3 (keeps prefill path working)
w_fused = torch.cat([w1_3d, w3_3d], dim=-1)\
            .reshape(args.dim, 2 * args.hidden_dim)

# mlp.py forward()
use_fused = self.fuse_gate_up and mode == Mode.DECODE
if use_fused:
    fused = ttnn.linear(x, self.w_gate_up, ...)
    fused = ttnn.to_memory_config(fused, ttnn.L1_MEMORY_CONFIG)
    w1_out = fused[..., :H]    # slice
    w3_out = fused[..., H:]    # slice
    # then existing ttnn.mul with activations=[SILU]
\end{lstlisting}
  \vfill
  \small
  Guarded on DECODE so prefill, Galaxy, and Llama-prefetcher paths are
  untouched. Took two passes to get right --- prefill
  \texttt{binary\_ng} barfed the first time.
\end{frame}

\begin{frame}{Result (on top of HiFi2)}
  \centering
  \small
  \begin{tabular}{lrrr}
    \toprule
     & HiFi2 only & +Fusion & $\Delta$ \\
    \midrule
    Matmul total (ms)        & 10.45 & 10.28 & ${-}1.6\%$ \\
    \quad calls               &  180 &   150 &  \\
    BinaryNg (ms)            &  0.69 & \alert{0.29} & \alert{${-}58\%$} \\
    Slice (new) + reshard    &  0.08 &  0.24 & $+0.16$ \\
    \midrule
    \emph{Total device (ms)} & \emph{36.75} & \alert{\emph{36.38}} & \alert{${-}1.0\%$} \\
    \bottomrule
  \end{tabular}
  \vfill
  \textbf{The matmul gain was smaller than we predicted.} The doubled core
  grid helps, but decode batch-1 matmul is DRAM-bandwidth-bound on the
  weight read --- more cores wait on the same DRAM bus.\\[0.5em]
  The \emph{real} win showed up somewhere we didn't expect: the
  downstream \texttt{binary\_ng} eltwise got 58\% cheaper because it
  now consumes L1-interleaved instead of L1-width-sharded inputs.
\end{frame}

% ---------------------------------------------------------------------------
\section{End-to-End}

\begin{frame}{All three attempts, side by side}
  \centering
  \small
  \begin{tabular}{lrrr}
    \toprule
    Configuration                       & Total (ms) & Std & vs.\ baseline \\
    \midrule
    Baseline                            & 41.78 & 0.008 & --- \\
    $+$ HiFi2                           & 36.78 & 0.009 & ${-}12.0\%$ \\
    $+$ HiFi2 $+$ Fusion                & \textbf{36.38} & 0.005 & \textbf{${-}12.9\%$} \\
    \midrule
    TopK padding (failed)               & \alert{56.80} & 0.032 & \alert{$+36.0\%$} \\
    \bottomrule
  \end{tabular}
  \vfill
  \begin{itemize}
    \item Numbers are mean of 5 repeated runs. Std is in
          \emph{microseconds}; deltas are in milliseconds. Signal-to-noise
          ratio 150:1 to 3{,}000:1 — these are not measurement noise.
    \item Per-iter decode: $10.4~\text{ms} \to 9.1~\text{ms}$. Throughput at
          10 layers: $\sim$96 $\to$ $\sim$110 tok/s/user.
    \item \texttt{ci-token-matching} passes. No generation-quality regression.
  \end{itemize}
\end{frame}

% ---------------------------------------------------------------------------
\begin{frame}{Does the speedup hold at other context lengths?}
  \small
  \centering
  \begin{tabular}{lrrr}
    \toprule
     & seq=512 & seq=1024 & seq=2048 \\
    \midrule
    Baseline               & 41.78 & 41.77 & 41.78 \\
    $+$ HiFi2              & 36.76 (${-}12.00\%$) & 36.77 (${-}11.98\%$) & 36.78 (${-}11.96\%$) \\
    $+$ HiFi2 $+$ Fusion   & \textbf{36.40} (${-}12.88\%$) & \textbf{36.39} (${-}12.88\%$) & \textbf{36.41} (${-}12.85\%$) \\
    \bottomrule
  \end{tabular}
  \vfill
  \begin{itemize}
    \item Speedup is \textbf{flat across input shapes} (variation $<$ 0.05pp).
    \item Absolute baseline is also flat — at decode batch 1 the cost is
          dominated by DRAM weight reads and the 1-core TopK; neither scales
          with cached-context length.
    \item SDPA does grow with seq\_len, but starts at 0.4\% of total, so
          the growth disappears under the microsecond noise floor.
  \end{itemize}
\end{frame}

% ---------------------------------------------------------------------------
\section{Lessons}

\begin{frame}{What we'd tell someone starting out}
  \begin{enumerate}
    \item \textbf{Profile before optimizing.} Our biggest ``fix'' turned
          out to be against a one-core sampling op nobody was looking at.
          Without Tracy we'd have spent the semester tuning attention.
    \item \textbf{On \texttt{tt-metal}, shape is not dispatch.} Padding
          inputs to look friendlier is a CUDA reflex. Here the
          \emph{device operation} class chooses the kernel, and shape is
          only one input to that decision.
    \item \textbf{Expected wins are often replaced by unexpected ones.}
          Fusion's matmul payoff was small. The payoff that actually
          mattered was a cheaper eltwise on interleaved operands.
          Re-profile after every change.
    \item \textbf{Negative results are worth writing up.} Our TopK
          padding attempt is something somebody else was eventually
          going to try. Now they don't have to.
  \end{enumerate}
\end{frame}

\begin{frame}{Future work}
  \begin{itemize}
    \item \textbf{Parallelize TopK.} Largest single number in the
          profile. C++ kernel work.
    \item \textbf{DRAM weight prefetcher for Qwen.} Exists for
          Llama-3.1-8B. Would overlap weight fetch with compute for
          the MLP matmuls and help the bandwidth-bound fused matmul we
          built.
    \item \textbf{Async AllGather $+$ matmul.} 4.8\% CCL overhead still
          sits in the critical path. Reference implementation:
          \texttt{experimental/ccl/llama\_all\_gather\_matmul\_async/}.
  \end{itemize}
\end{frame}

% ---------------------------------------------------------------------------
\begin{frame}[standout]
  Thank you.\\[0.5em]
  Questions?
\end{frame}

% ---------------------------------------------------------------------------
\appendix

\begin{frame}[fragile]{Backup: Reproducing a profile}
\begin{lstlisting}
cd /home/rassulmagauin/tt-metal
source python_env/bin/activate
export HF_MODEL=/home/rassulmagauin/models/Qwen2.5-7B-Instruct
export TT_METAL_HOME=/home/rassulmagauin/tt-metal
export MESH_DEVICE=P150x4

python3 -m tracy -p -r \
  -o generated/profiler/reports/<name> \
  --check-exit-code -a device_kernel_duration -t 5000 \
  -m 'pytest models/tt_transformers/demo/simple_text_demo.py \
      -k "device-perf and performance" \
      --num_layers 10 --batch_size 1 --max_seq_len 1024 \
      --max_generated_tokens 2 --mode decode --paged_attention 1 -x'
\end{lstlisting}
  \vfill
  Analyze with \texttt{tt-perf-report} on the emitted
  \texttt{ops\_perf\_results\_*.csv}. $\sim$90~s wall time per run on a
  warm cache.
\end{frame}

\begin{frame}{Backup: Where the code lives}
  \small
  Changes are on top of tt-metal commit \texttt{a4d8480d3e}. Three
  commits, two of them shipped:
  \vfill
  \begin{itemize}
    \item \texttt{965e0f4622} --- HiFi4 $\to$ HiFi2 for MLP FF1/FF2/FF3
          in \texttt{performance\_decoder\_config.json} (all 28 layers).
    \item \texttt{e05a044c4f} --- Gate/up projection fusion in
          \texttt{models/tt\_transformers/tt/mlp.py}. Gated on
          decode mode; prefill and Galaxy configs use the original
          two-linear path.
    \item \texttt{f6de95bb02} --- Revert of the TopK vocab padding in
          \texttt{model\_config.py}. Kept for provenance.
  \end{itemize}
\end{frame}

\begin{frame}{Backup: Compile-cache variance}
  Back-to-back runs of architecturally identical models showed
  $\sim$44\% variation in matmul mean time per op, traced to
  compile-cache state.
  \vfill
  For this report, all three comparison runs (Tables 2, 3, 4 of the
  paper) were taken with warm caches inside the same shell session.
  Deltas are stable under repeated measurement.
  \vfill
  Raw CSVs are under\\
  \small
  \texttt{generated/profiler/reports/}\\
  \texttt{qwen25\_7b\_decode\_*/}\\
  in the \texttt{tt-metal} working tree.
\end{frame}

\end{document}
